\section{Results with stronger dropout}
\label{sec:dropoutresults}
Appendix Figure~\ref{fig:dropout} compares validation loss with dropout 0.1 and 0.2 for each model. The full follow-up curves are also available in the repository’s regularization analysis notebook.

\begin{table}[H]
\centering
\caption{Effect of dropout on sampled validation losses. $\Delta$ best is dropout 0.2 minus 0.1; negative favors the intervention.}
\begin{tabular}{lrrrrr}
\toprule & \multicolumn{2}{c}{Best validation} & & \multicolumn{2}{c}{Final validation}\\
Model & $p=0.1$ & $p=0.2$ & $\Delta$ best & $p=0.1$ & $p=0.2$\\\midrule
Small & \textbf{1.7652} & 1.7665 & $+0.0013$ & 1.7659 & 1.7665\\
Medium & 1.8018 & \textbf{1.7872} & $-0.0146$ & 1.8892 & 1.8051\\
Large & 1.8122 & \textbf{1.7968} & $-0.0154$ & 1.9224 & 1.8305\\\bottomrule
\end{tabular}
\end{table}

\subsection*{How the benefit differs by model size}
With stronger dropout, medium and large reduce their best validation perplexities by about 1.45\% and 1.53\%. Medium reaches its best recorded result at 60.1 minutes instead of 40.1, while large reaches it at 60.4 minutes instead of 50.3. Their final validation losses improve by 0.0840 and 0.0919 nats/token. These improvements are larger than the changes in their best losses. Stronger dropout therefore mainly reduces the worsening seen late in training, while also slightly improving the best checkpoints.

The final gaps between training and validation loss fall from 0.5593 to 0.3841 for medium and from 0.6360 to 0.4572 for large. Validation loss still rises after its minimum, so dropout reduces overfitting without removing it. Since we evaluate about every ten minutes, we cannot identify the exact best stopping time.

Small's best loss increases by just $0.0013$, or about 0.13\% in perplexity. Its final training--validation gap falls from 0.2055 to 0.1444, but validation performance does not improve. This shows why a smaller gap alone does not mean a better model. The recorded early validation curves are worse with stronger dropout for all three sizes, but it appears more useful later for the larger models. The curves support this interpretation, but do not prove which internal features caused the result.

\subsection*{What remains unanswered before testing}
Stronger dropout helps medium and large, but neither overtakes small. This suggests that overfitting explains part of their later worsening, while stronger regularization alone does not resolve the size ranking. Small is an unsuccessful intervention: its smaller gap does not produce better validation predictions. Its $0.0013$ loss difference is too small to interpret confidently without repeated seeds. We should not describe dropout 0.2 as harmful in general.

The next questions are whether these small improvements repeat across seeds, whether another dropout or weight-decay value helps more, and whether the ranking changes with a different learning rate or equal-token budget. These remain open because those experiments were not run. Testing now addresses a narrower question: how do the validation-selected settings predict unseen pieces under one common scoring method?

\subsection*{How comparable the runs are}
The follow-up runs process 4.139B, 1.329B, and 0.507B tokens, close to the baseline counts. A large change in training speed is therefore unlikely to explain the pattern, although equal time does not guarantee equal amounts of training. Small still has the lowest sampled validation loss across both dropout settings. These runs suggest that the benefit of dropout depends on model size; they do not show that 0.2 is always best or better than tuning weight decay. The six training runs use about 12 GPU-hours of loop time, excluding benchmarking, setup, and final evaluation.
